% -------------------
%    bibliography
% -------------------
\usepackage[backend=biber,style=gost-numeric,sorting=none]{biblatex}
\usepackage{etoolbox}

% Запрет переноса между инициалами и фамилией ("А. Ф.~Стефаниди"),
% а также между двумя инициалами. Норма русской научной верстки.
\setcounter{highnamepenalty}{10000}
\setcounter{lownamepenalty}{10000}

\AtBeginBibliography{\hypersetup{hidelinks}}
\DeclareFieldFormat{author}{\textnormal{#1}}
\DeclareFieldFormat{editor}{\textnormal{#1}}
\renewbibmacro*{author}{%
  \printnames[author]{author}%
}

% ГОСТ 7.32 п.4.10.3: нумеровать арабскими цифрами с точкой и печатать с
% абзацного отступа. Номер «N.» стоит с абзацного отступа (1.25 см) —
% как обычный абзац, — текст идёт сразу за номером, переносы на новую
% строку выравниваются по левому полю (висячий отступ).
\defbibenvironment{bibliography}
  {\list
     {\printfield{labelnumber}.}              % метка вида «1.»
     {\setlength{\labelwidth}{0pt}%
      \setlength{\labelsep}{0.4em}%           % один пробел между «1.» и текстом
      \setlength{\leftmargin}{0pt}%           % перенос — по левому полю
      \setlength{\itemindent}{1.25cm}%        % первая строка — с абзацного отступа
      \setlength{\itemsep}{0pt}%
      \setlength{\parsep}{0pt}%
      \setlength{\topsep}{0pt}%
      \setlength{\partopsep}{0pt}%
      \setlength{\listparindent}{1.25cm}%
      \setlength{\rightmargin}{0pt}%
      \setlength{\parskip}{0pt}%
      \setlength{\parindent}{0pt}%
      \renewcommand{\makelabel}[1]{##1\hfil}% выравнивание метки по левому краю
     }}
  {\endlist}
  {\item}

\renewrobustcmd*{\mkbibemph}[1]{\normalfont#1}
\ExecuteBibliographyOptions{maxbibnames=99, minbibnames=1}
\addbibresource{lib/biblio.bib}
\renewcommand*{\bibfont}{\setlength{\parindent}{1.25cm}%
\setmainfont{Times New Roman}\fontsize{14pt}{16pt}\selectfont}
\usepackage{booktabs}
\usepackage{pdflscape}
\DefineBibliographyStrings{russian}{
    urlseen = {дата обращения} % Текст перед датой обращения
}

% Переопределяем, как выводится URL и дата обращения
% Это более сложная настройка, если стандартный вывод gost-numeric не устраивает.
% Сначала попробуйте просто с urldate и note/[howpublished] в .bib файле.
% Если нужно тонко настроить, то можно так:
\DeclareFieldFormat[online]{url}{%
  {[Электронный ресурс]. URL:\space\nolinkurl{#1}}% <--- \nolinkurl убирает кликабельность
}
% или если всегда добавлять [Электронный ресурс] перед URL для @online

\DeclareFieldFormat{urldate}{\mkbibparens{\bibstring{urlseen}\addcolon\space#1}}

% Чтобы "URL:" появлялось перед самой ссылкой
% (стили ГОСТ иногда требуют это, иногда нет)
\DefineBibliographyStrings{russian}{
  url = {URL} % Меняем "Режим доступа" на "URL"
}

\DeclareFieldFormat[article,inproceedings,incollection,book,thesis,report,patent,online,misc,unpublished]{title}{\normalfont{#1}} % Явно для основных типов
\DeclareFieldFormat{booktitle}{\normalfont{#1}}    % Для названия сборника или трудов конференции
\DeclareFieldFormat{journaltitle}{\normalfont{#1}} % Для названия журнала
\DeclareFieldFormat{issuetitle}{\normalfont{#1}}   % Для названия выпуска журнала
\DeclareFieldFormat{eventtitle}{\normalfont{#1}}   % Для названия мероприятия (конференции)
\setlength{\bibindent}{1.25cm} % Отступ слева для каждого элемента списка
